\mainchapter{Resultados}{\topbanner}\label{cap:resultados}
\vspace*{0.7cm}

En esta sección se presentan los resultados obtenidos a lo largo de las diferentes etapas del proyecto. Primero se exponen los resultados de la identificación y selección de la tecnología, incluyendo el análisis DAR que condujo a la elección de FreeIPA. Luego se describen los resultados del diseño de la arquitectura y de la implementación del servicio de directorio, incluyendo la integración con los servicios contemplados en el alcance. A continuación se presentan los resultados de las pruebas de verificación y el grado de cumplimiento de los requisitos establecidos. Finalmente, se ofrece una síntesis consolidada de los principales hallazgos del proyecto.

\section{Resultados de la identificación y selección de la tecnología}
\label{sec:resultados-seleccion}

En esta sección se presentan los resultados obtenidos durante el proceso de identificación y selección de la tecnología para la implementación del servicio de directorio del Grupo de Investigación en Redes, Información y Distribución (GRID). A partir de las necesidades identificadas, los criterios de evaluación definidos y el análisis DAR realizado, se determinó la tecnología que presentó mayor correspondencia con los requerimientos establecidos para la solución.

\subsection{Resultado del análisis DAR}
\label{subsec:resultado-dar}

El análisis DAR estructuró la comparación de las tecnologías consideradas para la implementación del servicio de directorio, tomando como referencia los criterios definidos durante la fase de identificación de tecnologías. Como resultado de esta evaluación, se seleccionó FreeIPA como la tecnología base para la implementación de la solución.

La selección se fundamentó en su capacidad para proporcionar un servicio centralizado de gestión de identidades, administración de usuarios y grupos, aplicación de políticas de seguridad y mecanismos de autenticación mediante protocolos compatibles con los servicios considerados en el proyecto. Asimismo, la solución incorpora una arquitectura con capacidad de replicación, lo que contribuye a la disponibilidad del servicio de directorio.

\section{Resultados del diseño de la solución}
\label{sec:resultados-diseno}

Como resultado de la fase de diseño se obtuvo una arquitectura de servicio de directorio orientada a centralizar la gestión de identidades y el control de acceso de los servicios considerados en el proyecto. El diseño contempló la utilización de un servidor FreeIPA principal y un servidor réplica, junto con un punto de acceso común para los servicios que requieren autenticación mediante LDAP.

La arquitectura resultante separa los componentes destinados a la gestión de identidades de los servicios consumidores del directorio. De esta manera, los usuarios y grupos son administrados desde el servicio de directorio, mientras que los servicios integrados utilizan esta información para realizar los procesos de autenticación y autorización correspondientes.

\section{Resultados de la implementación}
\label{sec:resultados-implementacion}

Durante la implementación se desplegó el servicio de directorio sobre FreeIPA y se configuraron los componentes requeridos para la gestión centralizada de identidades: cuentas de usuario, grupos y políticas de seguridad necesarias para las pruebas de integración con los servicios definidos en el alcance del proyecto.

Asimismo, fue implementado un segundo servidor FreeIPA como réplica del principal. Esta configuración mantiene una copia sincronizada de la información administrada por el servicio de directorio y crea las condiciones para evaluar la continuidad del servicio ante la indisponibilidad de uno de los servidores.

Con la configuración quedaron vigentes las políticas de contraseñas y los mecanismos de administración de usuarios y grupos definidos para la solución. Estos ajustes aplican las restricciones de seguridad establecidas durante la definición de requisitos y proveen una estructura de identidades utilizable por los servicios integrados.

La configuración de la réplica sincronizó los servidores FreeIPA desplegados. Con ello, la arquitectura gana redundancia del servicio de directorio y reduce la dependencia de una única instancia para la disponibilidad de la información de identidades.

Sobre esta misma configuración se ejecutaron las pruebas relacionadas con la disponibilidad y replicación de la información, verificando el comportamiento de la solución ante escenarios de indisponibilidad del servidor principal.

\subsection{Integración con los servicios}
\label{subsec:resultado-integracion}

La solución integró el servicio de directorio con los servicios contemplados en el proyecto mediante el protocolo LDAP. Para cada servicio se configuraron los parámetros necesarios para consultar el directorio, autenticar los usuarios y aplicar las restricciones de acceso asociadas a los grupos definidos.

Como resultado, se obtuvo un mecanismo centralizado mediante el cual las credenciales y pertenencia a grupos administradas en FreeIPA pueden ser utilizadas por los servicios integrados, evitando la necesidad de mantener de forma independiente las cuentas de usuario en cada servicio.

\section{Resultados de las pruebas}
\label{sec:resultados-pruebas}

Las pruebas verificaron el comportamiento de la solución implementada frente a los requisitos establecidos para el servicio de directorio. Su ejecución abarcó el servicio de directorio, su configuración de replicación y los servicios integrados, utilizando los procedimientos definidos durante la etapa de validación.

A partir de las evidencias obtenidas se determinó el grado de cumplimiento de cada requisito, diferenciando aquellos que fueron implementados satisfactoriamente de aquellos que presentaron limitaciones derivadas de la arquitectura seleccionada o del alcance establecido para el proyecto.

\subsection{Cumplimiento de requisitos}
\label{subsec:resultado-requisitos}

La evaluación de los requisitos determinó que la solución implementada satisface principalmente las necesidades relacionadas con la centralización de identidades, la administración de usuarios y grupos, la aplicación de políticas de seguridad, la integración mediante LDAP y la disponibilidad del servicio mediante replicación.

El requisito relacionado con el control y monitoreo de accesos se abordó mediante la utilización de grupos para determinar la autorización de los usuarios en los servicios integrados. De esta forma, la pertenencia de un usuario a los grupos definidos determina si este puede acceder al servicio correspondiente.

Por otra parte, el registro y trazabilidad de las operaciones relacionadas con el servicio de directorio se evaluó a partir de los mecanismos de registro proporcionados por FreeIPA y sus componentes. Estos mecanismos permiten disponer de información asociada a las operaciones realizadas sobre el servicio y constituyen un elemento de apoyo para las actividades de auditoría.

Los resultados también pusieron de manifiesto limitaciones en determinados requisitos. En particular, la autenticación multifactor no pudo ser incorporada directamente en las integraciones realizadas mediante LDAP, debido a que los mecanismos de autenticación utilizados por los servicios integrados no contemplan este factor dentro del flujo de autenticación implementado. De igual manera, la recuperación autónoma de contraseñas por parte de los usuarios no fue implementada, por lo que esta funcionalidad requiere mecanismos adicionales y no forma parte de la solución finalmente desplegada.

\subsection{Resultado global de la validación}
\label{subsec:resultado-global-validacion}

En términos generales, los resultados obtenidos evidencian que la solución implementada permite centralizar la gestión de identidades del entorno evaluado y proporcionar un mecanismo común de autenticación y autorización para los servicios integrados. La incorporación de un servidor réplica complementa esta capacidad al proporcionar redundancia para el servicio de directorio.

No obstante, la validación mostró igualmente que determinadas capacidades planteadas inicialmente requieren componentes o mecanismos adicionales a la integración LDAP utilizada en el proyecto. Estas limitaciones quedaron contempladas en la evaluación final de la solución y delimitan las funcionalidades efectivamente alcanzadas respecto de aquellas que podrían ser incorporadas en trabajos posteriores.

\section{Síntesis de los resultados}
\label{sec:sintesis-resultados}

Como resultado del proyecto se obtuvo una solución de servicio de directorio basada en FreeIPA, integrada con los servicios definidos en el alcance y respaldada por una arquitectura de replicación. La solución centraliza la gestión de identidades, usuarios y grupos, aplica políticas de seguridad y utiliza la información del directorio para controlar el acceso a los servicios integrados.

Las pruebas comprobaron las funcionalidades implementadas y determinaron el nivel de cumplimiento de los requisitos definidos. Asimismo, dejaron identificadas las limitaciones asociadas principalmente a las capacidades de autenticación multifactor y recuperación de contraseñas cuando la integración se realiza mediante LDAP. En consecuencia, los resultados obtenidos permiten valorar la solución implementada de acuerdo con las capacidades efectivamente alcanzadas y con las condiciones establecidas en el alcance del proyecto.
